\documentclass[11pt]{article}
\usepackage[utf8]{inputenc}
\usepackage{amsmath, amssymb, amsthm}
\usepackage{geometry}
\geometry{a4paper, margin=1in}
\usepackage{hyperref}
\usepackage{enumitem}

\title{Problem Sets: Undetermined Coefficients and the Annihilator Method}
\author{Based on Chapter 6}
\date{}

\begin{document}

\maketitle
\tableofcontents
\newpage

\section{Concept Checklist \& Question Mapping}
\label{sec:checklist_start}
The following checklist outlines the core concepts, problem types, and common traps found in the text. The corresponding problem numbers are listed next to each concept to ensure comprehensive coverage.

\begin{itemize}
    \item \textbf{C1: Finding Differential Annihilators (Basic \& Composite)} - Using the annihilator tables for polynomials, exponentials, sines/cosines, and their products/sums. \\
    \textit{Questions: 1, 2, 3, 4, 5}
    \item \textbf{C2: Determining the Form of $y_p$ (No Overlap)} - Setting up the particular solution guess without needing the "s" multiplier. \\
    \textit{Questions: 6}
    \item \textbf{C3: Determining the Form of $y_p$ (With Overlap)} - Identifying overlap with the complementary solution $y_c$ and applying the "s" multiplier trick. \\
    \textit{Questions: 7, 8, 9, 10}
    \item \textbf{C4: Using the Annihilator Method to Determine $y_p$ Form} - Applying the operator approach rigorously to justify the form of the guess. \\
    \textit{Questions: 11, 12, 13, 14, 15}
    \item \textbf{C5: Finding the General Solution (Full Solution)} - Executing the full method of undetermined coefficients to find the exact constants. \\
    \textit{Questions: 16, 17, 18, 19, 20}
    \item \textbf{C6: Initial Value Problems (IVPs)} - Applying initial conditions to the \textit{full} general solution $y(x) = y_c(x) + y_p(x)$. \\
    \textit{Questions: 21, 22, 23}
    \item \textbf{C7: Systems of Differential Equations} - Using operator notation and the elimination method to decouple and solve systems. \\
    \textit{Questions: 24, 25}
    \item \textbf{C8: Identifying Method Limitations} - Recognizing when the annihilator method fails (variable coefficients, non-table forcing functions like secant/tangent). \\
    \textit{Questions: 26, 27}
    \item \textbf{C9: "Find the Flaw" / Conceptual Traps} - Diagnosing common errors such as the Overlap Trap, Incomplete Guess Trap, and Annihilator of a Sum vs. Product. \\
    \textit{Questions: 28, 29}
    \item \textbf{C10: Proofs and Derivations} - Proving annihilator properties using operator algebra. \\
    \textit{Questions: 30}
\end{itemize}

\newpage

\section{Problem Sets}

\subsection*{Part 1: Finding Differential Annihilators}
Find a differential operator with constant coefficients that annihilates the following functions:
\begin{enumerate}
    \item $f(x) = x^3 - 2x + 5$
    \item $f(x) = e^{-4x}$
    \item $f(x) = x e^{2x} \cos(3x)$
    \item $f(x) = 4e^{x} + \sin(2x)$
    \item $f(x) = x^2 e^{-x} \sin(x) + x^5$
\end{enumerate}

\subsection*{Part 2: Determining the Form of a Particular Solution (Setup Only)}
Determine the form of a particular solution $y_p(x)$ for the following differential equations using the Method of Undetermined Coefficients. Do not solve for the constants.
\begin{enumerate}[resume]
    \item $y'' - 5y' + 6y = e^{4x} + x$
    \item $y'' - 4y' + 4y = e^{2x}$
    \item $y''' + y' = \sin x + x^2$
    \item $y'' - y = x e^x$
    \item $y^{(4)} - y = e^x + \cos x$
\end{enumerate}

\subsection*{Part 3: Using the Annihilator Method to Determine the Form of $y_p$}
Use the Annihilator Method (operator approach) to rigorously determine the form of the particular solution $y_p(x)$ for the following equations.
\begin{enumerate}[resume]
    \item $y'' - 3y' - 4y = e^{-x}$
    \item $y'' + 4y = \cos(2x)$
    \item $y'' - 2y' + y = x^2$
    \item $y''' - y'' = 6$
    \item $y'' - y = x e^{-x}$
\end{enumerate}

\subsection*{Part 4: Finding the General Solution}
Find the general solution to the following nonhomogeneous differential equations.
\begin{enumerate}[resume]
    \item $y'' - 5y' + 4y = 8e^{3x}$
    \item $y'' + y = 2 \sin x$
    \item $y'' - 2y' + y = 3e^x$
    \item $y'' - y = x^2$
    \item $y''' + y'' = 2x$
\end{enumerate}

\subsection*{Part 5: Initial Value Problems}
Solve the following initial value problems.
\begin{enumerate}[resume]
    \item $y'' + 4y = 8; \quad y(0) = 2, \quad y'(0) = 0$
    \item $y'' - 3y' + 2y = 2e^{3x}; \quad y(0) = 1, \quad y'(0) = 2$
    \item $y'' - y = 2x; \quad y(0) = 0, \quad y'(0) = 0$
\end{enumerate}

\subsection*{Part 6: Systems of Differential Equations (Elimination Method)}
Use the elimination method to find a general solution to the following systems.
\begin{enumerate}[resume]
    \item $\frac{dx}{dt} + y = t$ \\ $\frac{dy}{dt} - x = 0$
    \item $\frac{dx}{dt} - 2x + \frac{dy}{dt} = e^t$ \\ $\frac{dx}{dt} + \frac{dy}{dt} - y = 0$
\end{enumerate}

\subsection*{Part 7: Conceptual Variations and Limitations}
\begin{enumerate}[resume]
    \item Explain why the Method of Undetermined Coefficients / Annihilator Method cannot be used to solve $y'' + y = \tan x$. What method must be used instead?
    \item Explain why the Annihilator Method cannot be used to solve $x^2 y'' - 4x y' + 6y = x^3$.
\end{enumerate}

\subsection*{Part 8: Advanced Diagnostic Testing: "Find the Flaw"}
Identify the critical flaw in the following proposed solutions, explain why it is wrong, and state the correct next step.
\begin{enumerate}[resume]
    \item \textbf{Problem:} Solve $y'' - 9y = e^{3x}$. \\
    \textbf{Flawed Solution:} The roots of the characteristic equation are $r = \pm 3$. The forcing function is $e^{3x}$, so we guess $y_p = A e^{3x}$. Substituting this into the ODE gives $9A e^{3x} - 9A e^{3x} = e^{3x}$, which simplifies to $0 = e^{3x}$. Since this is impossible, there is no particular solution.
    \item \textbf{Problem:} Solve $y'' + 4y = x^2$. \\
    \textbf{Flawed Solution:} The forcing function is $x^2$, which is a polynomial. We guess $y_p = A x^2$. The second derivative is $y_p'' = 2A$. Substituting into the ODE gives $2A + 4A x^2 = x^2$. Equating coefficients gives $4A = 1 \implies A = 1/4$, and $2A = 0 \implies A = 0$. This is a contradiction.
\end{enumerate}

\subsection*{Part 9: Proofs and Derivations}
\begin{enumerate}[resume]
    \item Use the properties of linear differential operators to prove that the annihilator for $f(x) = x e^{\alpha x}$ is $(D - \alpha)^2$.
\end{enumerate}

\newpage

\section{Detailed Solutions}

\subsection*{Part 1: Finding Differential Annihilators}

\textbf{Solution 1:} $f(x) = x^3 - 2x + 5$. \\
This is a polynomial of degree 3. The $(m+1)$-th derivative of a polynomial of degree $m$ is zero. Therefore, the annihilator is $D^{3+1} = D^4$.

\textbf{Solution 2:} $f(x) = e^{-4x}$. \\
For an exponential function $e^{\alpha x}$, the annihilator is $(D - \alpha)$. Here, $\alpha = -4$. The annihilator is $D - (-4) = D + 4$.

\textbf{Solution 3:} $f(x) = x e^{2x} \cos(3x)$. \\
This is of the form $x^k e^{\alpha x} \cos(\beta x)$ with $k=1$, $\alpha=2$, and $\beta=3$. The base annihilator for $e^{2x} \cos(3x)$ is $(D - \alpha)^2 + \beta^2 = (D - 2)^2 + 9$. Because it is multiplied by $x^1$, we raise the operator to the power of $k+1 = 2$. The annihilator is $[(D - 2)^2 + 9]^2$.

\textbf{Solution 4:} $f(x) = 4e^{x} + \sin(2x)$. \\
We find the annihilator for each term and multiply them (Annihilator of a sum is the product of annihilators). For $4e^x$, $\alpha=1$, so $A_1 = (D - 1)$. For $\sin(2x)$, $\alpha=0, \beta=2$, so $A_2 = (D^2 + 4)$. The total annihilator is $(D - 1)(D^2 + 4)$.

\textbf{Solution 5:} $f(x) = x^2 e^{-x} \sin(x) + x^5$. \\
Term 1: $x^2 e^{-x} \sin(x)$. Here $k=2, \alpha=-1, \beta=1$. Base operator: $(D+1)^2 + 1$. Raised to $k+1=3$: $[(D+1)^2 + 1]^3$. \\
Term 2: $x^5$. Polynomial of degree 5. Annihilator is $D^6$. \\
Total annihilator: $[(D+1)^2 + 1]^3 D^6$.

\subsection*{Part 2: Determining the Form of a Particular Solution (Setup Only)}

\textbf{Solution 6:} $y'' - 5y' + 6y = e^{4x} + x$. \\
Characteristic equation: $r^2 - 5r + 6 = (r-2)(r-3) = 0$. Roots: $r=2, 3$. \\
$y_c(x) = c_1 e^{2x} + c_2 e^{3x}$. \\
Forcing function $f(x) = e^{4x} + x$. \\
Base guess for $e^{4x}$ is $A e^{4x}$. No overlap with $y_c$. \\
Base guess for $x$ is $Bx + C$. No overlap with $y_c$. \\
Final form: $y_p(x) = A e^{4x} + Bx + C$.

\textbf{Solution 7:} $y'' - 4y' + 4y = e^{2x}$. \\
Characteristic equation: $r^2 - 4r + 4 = (r-2)^2 = 0$. Roots: $r=2$ (multiplicity 2). \\
$y_c(x) = c_1 e^{2x} + c_2 x e^{2x}$. \\
Forcing function $f(x) = e^{2x}$. Base guess is $A e^{2x}$. \\
Overlap check: $e^{2x}$ and $x e^{2x}$ are already in $y_c$. We must multiply by $x^2$. \\
Final form: $y_p(x) = A x^2 e^{2x}$.

\textbf{Solution 8:} $y''' + y' = \sin x + x^2$. \\
Characteristic equation: $r^3 + r = r(r^2 + 1) = 0$. Roots: $r=0, \pm i$. \\
$y_c(x) = c_1 + c_2 \cos x + c_3 \sin x$. \\
For $f_1(x) = \sin x$, base guess is $A \cos x + B \sin x$. This overlaps entirely with $y_c$. Multiply by $x$: $x(A \cos x + B \sin x)$. \\
For $f_2(x) = x^2$, base guess is $Cx^2 + Dx + E$. The constant term $E$ overlaps with $c_1$ in $y_c$. Multiply by $x$: $x(Cx^2 + Dx + E) = Cx^3 + Dx^2 + Ex$. \\
Final form: $y_p(x) = x(A \cos x + B \sin x) + Cx^3 + Dx^2 + Ex$.

\textbf{Solution 9:} $y'' - y = x e^x$. \\
Characteristic equation: $r^2 - 1 = 0$. Roots: $r = \pm 1$. \\
$y_c(x) = c_1 e^x + c_2 e^{-x}$. \\
For $f(x) = x e^x$, base guess is $(Ax + B)e^x = A x e^x + B e^x$. \\
The term $B e^x$ overlaps with $c_1 e^x$ in $y_c$. Multiply the entire guess by $x$. \\
Final form: $y_p(x) = x(Ax + B)e^x = (Ax^2 + Bx)e^x$.

\textbf{Solution 10:} $y^{(4)} - y = e^x + \cos x$. \\
Characteristic equation: $r^4 - 1 = (r^2-1)(r^2+1) = 0$. Roots: $r = \pm 1, \pm i$. \\
$y_c(x) = c_1 e^x + c_2 e^{-x} + c_3 \cos x + c_4 \sin x$. \\
For $f_1(x) = e^x$, base guess $A e^x$ overlaps with $c_1 e^x$. Multiply by $x$: $A x e^x$. \\
For $f_2(x) = \cos x$, base guess $B \cos x + C \sin x$ overlaps with $c_3 \cos x + c_4 \sin x$. Multiply by $x$: $x(B \cos x + C \sin x)$. \\
Final form: $y_p(x) = A x e^x + x(B \cos x + C \sin x)$.

\subsection*{Part 3: Using the Annihilator Method to Determine the Form of $y_p$}

\textbf{Solution 11:} $y'' - 3y' - 4y = e^{-x}$. \\
Operator form: $L[y] = (D^2 - 3D - 4)y = (D-4)(D+1)y = e^{-x}$. \\
Annihilator for $e^{-x}$ is $A = (D+1)$. Apply $A$ to both sides: \\
$(D+1)(D-4)(D+1)y = 0 \implies (D+1)^2(D-4)y = 0$. \\
Roots of new equation: $r = 4, -1, -1$. \\
General solution to annihilated equation: $y(x) = c_1 e^{4x} + c_2 e^{-x} + c_3 x e^{-x}$. \\
The terms $c_1 e^{4x} + c_2 e^{-x}$ form $y_c(x)$. The remaining term gives the form of $y_p$: \\
$y_p(x) = A x e^{-x}$.

\textbf{Solution 12:} $y'' + 4y = \cos(2x)$. \\
Operator form: $L[y] = (D^2 + 4)y = \cos(2x)$. \\
Annihilator for $\cos(2x)$ is $A = (D^2 + 4)$. Apply $A$ to both sides: \\
$(D^2 + 4)(D^2 + 4)y = 0 \implies (D^2 + 4)^2 y = 0$. \\
Roots: $r = \pm 2i$ (multiplicity 2). \\
$y(x) = c_1 \cos(2x) + c_2 \sin(2x) + c_3 x \cos(2x) + c_4 x \sin(2x)$. \\
The terms $c_1 \cos(2x) + c_2 \sin(2x)$ form $y_c(x)$. The remaining terms give $y_p$: \\
$y_p(x) = x(A \cos(2x) + B \sin(2x))$.

\textbf{Solution 13:} $y'' - 2y' + y = x^2$. \\
Operator form: $L[y] = (D-1)^2 y = x^2$. \\
Annihilator for $x^2$ is $A = D^3$. Apply $A$ to both sides: \\
$D^3(D-1)^2 y = 0$. \\
Roots: $r = 1, 1, 0, 0, 0$. \\
$y(x) = c_1 e^x + c_2 x e^x + c_3 + c_4 x + c_5 x^2$. \\
The terms $c_1 e^x + c_2 x e^x$ form $y_c(x)$. The remaining terms give $y_p$: \\
$y_p(x) = A x^2 + B x + C$.

\textbf{Solution 14:} $y''' - y'' = 6$. \\
Operator form: $L[y] = D^2(D-1)y = 6$. \\
Annihilator for $6$ (a constant) is $A = D$. Apply $A$ to both sides: \\
$D \cdot D^2(D-1)y = 0 \implies D^3(D-1)y = 0$. \\
Roots: $r = 1, 0, 0, 0$. \\
$y(x) = c_1 e^x + c_2 + c_3 x + c_4 x^2$. \\
The terms $c_1 e^x + c_2 + c_3 x$ form $y_c(x)$ (since $r=0$ had multiplicity 2 in $L$). The remaining term gives $y_p$: \\
$y_p(x) = A x^2$.

\textbf{Solution 15:} $y'' - y = x e^{-x}$. \\
Operator form: $L[y] = (D-1)(D+1)y = x e^{-x}$. \\
Annihilator for $x e^{-x}$ is $A = (D+1)^2$. Apply $A$ to both sides: \\
$(D+1)^2(D-1)(D+1)y = 0 \implies (D-1)(D+1)^3 y = 0$. \\
Roots: $r = 1, -1, -1, -1$. \\
$y(x) = c_1 e^x + c_2 e^{-x} + c_3 x e^{-x} + c_4 x^2 e^{-x}$. \\
The terms $c_1 e^x + c_2 e^{-x}$ form $y_c(x)$. The remaining terms give $y_p$: \\
$y_p(x) = A x^2 e^{-x} + B x e^{-x} = (Ax^2 + Bx)e^{-x}$.

\subsection*{Part 4: Finding the General Solution}

\textbf{Solution 16:} $y'' - 5y' + 4y = 8e^{3x}$. \\
Step 1: $y_c(x)$. $r^2 - 5r + 4 = (r-1)(r-4) = 0 \implies r=1, 4$. $y_c(x) = c_1 e^x + c_2 e^{4x}$. \\
Step 2: $y_p(x)$. Guess $y_p = A e^{3x}$. (No overlap). \\
$y_p' = 3A e^{3x}$, $y_p'' = 9A e^{3x}$. \\
Substitute: $9A e^{3x} - 5(3A e^{3x}) + 4(A e^{3x}) = 8e^{3x}$. \\
$(9 - 15 + 4)A = 8 \implies -2A = 8 \implies A = -4$. \\
General Solution: $y(x) = c_1 e^x + c_2 e^{4x} - 4e^{3x}$.

\textbf{Solution 17:} $y'' + y = 2 \sin x$. \\
Step 1: $y_c(x)$. $r^2 + 1 = 0 \implies r = \pm i$. $y_c(x) = c_1 \cos x + c_2 \sin x$. \\
Step 2: $y_p(x)$. Base guess $A \cos x + B \sin x$ overlaps. Multiply by $x$: $y_p = x(A \cos x + B \sin x)$. \\
$y_p' = A \cos x + B \sin x + x(-A \sin x + B \cos x)$. \\
$y_p'' = -2A \sin x + 2B \cos x - x(A \cos x + B \sin x)$. \\
Substitute: $[-2A \sin x + 2B \cos x - x(A \cos x + B \sin x)] + [x(A \cos x + B \sin x)] = 2 \sin x$. \\
$-2A \sin x + 2B \cos x = 2 \sin x$. \\
Equate coefficients: $-2A = 2 \implies A = -1$. $2B = 0 \implies B = 0$. \\
General Solution: $y(x) = c_1 \cos x + c_2 \sin x - x \cos x$.

\textbf{Solution 18:} $y'' - 2y' + y = 3e^x$. \\
Step 1: $y_c(x)$. $r^2 - 2r + 1 = (r-1)^2 = 0 \implies r=1, 1$. $y_c(x) = c_1 e^x + c_2 x e^x$. \\
Step 2: $y_p(x)$. Base guess $A e^x$ overlaps. $A x e^x$ overlaps. Guess $y_p = A x^2 e^x$. \\
$y_p' = A(x^2 + 2x)e^x$. \\
$y_p'' = A(x^2 + 4x + 2)e^x$. \\
Substitute: $A e^x (x^2 + 4x + 2) - 2A e^x (x^2 + 2x) + A x^2 e^x = 3e^x$. \\
$A e^x [x^2 - 2x^2 + x^2 + 4x - 4x + 2] = 3e^x \implies 2A e^x = 3e^x \implies A = 3/2$. \\
General Solution: $y(x) = c_1 e^x + c_2 x e^x + \frac{3}{2} x^2 e^x$.

\textbf{Solution 19:} $y'' - y = x^2$. \\
Step 1: $y_c(x)$. $r^2 - 1 = 0 \implies r = \pm 1$. $y_c(x) = c_1 e^x + c_2 e^{-x}$. \\
Step 2: $y_p(x)$. Guess $y_p = A x^2 + B x + C$. \\
$y_p' = 2Ax + B$, $y_p'' = 2A$. \\
Substitute: $2A - (A x^2 + B x + C) = x^2$. \\
$-A x^2 - B x + (2A - C) = x^2$. \\
Equate coefficients: $-A = 1 \implies A = -1$. $-B = 0 \implies B = 0$. $2A - C = 0 \implies -2 - C = 0 \implies C = -2$. \\
General Solution: $y(x) = c_1 e^x + c_2 e^{-x} - x^2 - 2$.

\textbf{Solution 20:} $y''' + y'' = 2x$. \\
Step 1: $y_c(x)$. $r^3 + r^2 = r^2(r+1) = 0 \implies r = 0, 0, -1$. $y_c(x) = c_1 + c_2 x + c_3 e^{-x}$. \\
Step 2: $y_p(x)$. Base guess for $2x$ is $Ax + B$. This overlaps with $c_1 + c_2 x$. Multiply by $x^2$: $y_p = A x^3 + B x^2$. \\
$y_p' = 3A x^2 + 2B x$. \\
$y_p'' = 6A x + 2B$. \\
$y_p''' = 6A$. \\
Substitute: $6A + (6A x + 2B) = 2x \implies 6A x + (6A + 2B) = 2x$. \\
Equate coefficients: $6A = 2 \implies A = 1/3$. $6A + 2B = 0 \implies 2 + 2B = 0 \implies B = -1$. \\
General Solution: $y(x) = c_1 + c_2 x + c_3 e^{-x} + \frac{1}{3}x^3 - x^2$.

\subsection*{Part 5: Initial Value Problems}

\textbf{Solution 21:} $y'' + 4y = 8; \quad y(0) = 2, \quad y'(0) = 0$. \\
Step 1: General Solution. $r^2 + 4 = 0 \implies r = \pm 2i$. $y_c = c_1 \cos(2x) + c_2 \sin(2x)$. \\
Guess $y_p = A$. $y_p'' = 0$. Substitute: $0 + 4A = 8 \implies A = 2$. \\
$y(x) = c_1 \cos(2x) + c_2 \sin(2x) + 2$. \\
Step 2: Apply IVs. $y(0) = c_1(1) + c_2(0) + 2 = 2 \implies c_1 + 2 = 2 \implies c_1 = 0$. \\
$y'(x) = -2c_1 \sin(2x) + 2c_2 \cos(2x)$. \\
$y'(0) = 2c_2 = 0 \implies c_2 = 0$. \\
Final Solution: $y(x) = 2$.

\textbf{Solution 22:} $y'' - 3y' + 2y = 2e^{3x}; \quad y(0) = 1, \quad y'(0) = 2$. \\
Step 1: General Solution. $r^2 - 3r + 2 = (r-1)(r-2) = 0 \implies r=1, 2$. $y_c = c_1 e^x + c_2 e^{2x}$. \\
Guess $y_p = A e^{3x}$. $y_p' = 3A e^{3x}, y_p'' = 9A e^{3x}$. \\
Substitute: $9A e^{3x} - 9A e^{3x} + 2A e^{3x} = 2e^{3x} \implies 2A = 2 \implies A = 1$. \\
$y(x) = c_1 e^x + c_2 e^{2x} + e^{3x}$. \\
Step 2: Apply IVs. $y(0) = c_1 + c_2 + 1 = 1 \implies c_1 + c_2 = 0$. \\
$y'(x) = c_1 e^x + 2c_2 e^{2x} + 3e^{3x}$. \\
$y'(0) = c_1 + 2c_2 + 3 = 2 \implies c_1 + 2c_2 = -1$. \\
System: 1) $c_1 + c_2 = 0$, 2) $c_1 + 2c_2 = -1$. Subtract (1) from (2): $c_2 = -1$. From (1), $c_1 = 1$. \\
Final Solution: $y(x) = e^x - e^{2x} + e^{3x}$.

\textbf{Solution 23:} $y'' - y = 2x; \quad y(0) = 0, \quad y'(0) = 0$. \\
Step 1: General Solution. $r^2 - 1 = 0 \implies r = \pm 1$. $y_c = c_1 e^x + c_2 e^{-x}$. \\
Guess $y_p = Ax + B$. $y_p'' = 0$. Substitute: $0 - (Ax + B) = 2x \implies -A = 2 \implies A = -2$, $-B = 0 \implies B = 0$. \\
$y(x) = c_1 e^x + c_2 e^{-x} - 2x$. \\
Step 2: Apply IVs. $y(0) = c_1 + c_2 - 0 = 0 \implies c_1 + c_2 = 0$. \\
$y'(x) = c_1 e^x - c_2 e^{-x} - 2$. \\
$y'(0) = c_1 - c_2 - 2 = 0 \implies c_1 - c_2 = 2$. \\
System: 1) $c_1 + c_2 = 0$, 2) $c_1 - c_2 = 2$. Add them: $2c_1 = 2 \implies c_1 = 1$. Then $c_2 = -1$. \\
Final Solution: $y(x) = e^x - e^{-x} - 2x$.

\subsection*{Part 6: Systems of Differential Equations (Elimination Method)}

\textbf{Solution 24:} \\
1) $Dx + y = t \implies y = t - Dx$ \\
2) $Dy - x = 0$ \\
Substitute $y$ into (2): $D(t - Dx) - x = 0 \implies 1 - D^2x - x = 0 \implies (D^2 + 1)x = 1$. \\
Solve for $x$: Characteristic equation $r^2 + 1 = 0 \implies r = \pm i$. $x_c(t) = c_1 \cos t + c_2 \sin t$. \\
Guess $x_p(t) = A$. $D^2(A) + A = 1 \implies A = 1$. \\
$x(t) = c_1 \cos t + c_2 \sin t + 1$. \\
Now find $y(t)$ using $y = t - Dx = t - x'(t)$: \\
$x'(t) = -c_1 \sin t + c_2 \cos t$. \\
$y(t) = t - (-c_1 \sin t + c_2 \cos t) = c_1 \sin t - c_2 \cos t + t$.

\textbf{Solution 25:} \\
1) $(D-2)x + Dy = e^t$ \\
2) $Dx + (D-1)y = 0$ \\
Multiply (1) by $D$ and (2) by $(D-2)$: \\
$D(D-2)x + D^2y = e^t$ (since $D(e^t) = e^t$) \\
$(D-2)Dx + (D-2)(D-1)y = 0$ \\
Subtract the second from the first: \\
$[D^2 - (D^2 - 3D + 2)]y = e^t \implies (3D - 2)y = e^t$. \\
Solve for $y$: $3y' - 2y = e^t$. Characteristic root $3r - 2 = 0 \implies r = 2/3$. $y_c(t) = c_1 e^{2t/3}$. \\
Guess $y_p(t) = A e^t$. Substitute: $3A e^t - 2A e^t = e^t \implies A = 1$. \\
$y(t) = c_1 e^{2t/3} + e^t$. \\
To find $x(t)$, use the original equations. From (1) minus (2): \\
$(-2x + y) = e^t \implies 2x = y - e^t \implies x(t) = \frac{1}{2}y(t) - \frac{1}{2}e^t$. \\
$x(t) = \frac{1}{2}(c_1 e^{2t/3} + e^t) - \frac{1}{2}e^t = \frac{1}{2}c_1 e^{2t/3}$.

\subsection*{Part 7: Conceptual Variations and Limitations}

\textbf{Solution 26:} \\
The Annihilator Method relies on finding a linear differential operator with constant coefficients that turns the forcing function $f(x)$ into zero. Functions like $\tan x$, $\sec x$, and $\ln x$ have infinitely many linearly independent derivatives and cannot be annihilated by any finite-order constant-coefficient operator. Therefore, the method fails. One must use the \textbf{Method of Variation of Parameters}.

\textbf{Solution 27:} \\
The Annihilator Method (and Undetermined Coefficients) strictly requires the differential equation to have \textbf{constant coefficients}. The equation $x^2 y'' - 4x y' + 6y = x^3$ has variable coefficients (it is a Cauchy-Euler equation). The operators $L$ and $A$ would not commute ($A[L[y]] \neq L[A[y]]$), breaking the fundamental algebraic justification of the method.

\subsection*{Part 8: Advanced Diagnostic Testing: "Find the Flaw"}

\textbf{Solution 28:} \\
\textbf{Flaw:} The student fell into the \textbf{Overlap Trap}. They failed to check their guess against the complementary solution $y_c(x)$. \\
\textbf{Correction:} The roots are $r = \pm 3$, so $y_c(x) = c_1 e^{3x} + c_2 e^{-3x}$. The forcing function $e^{3x}$ is already in $y_c(x)$. The correct next step is to multiply the guess by $x$, making it $y_p(x) = A x e^{3x}$.

\textbf{Solution 29:} \\
\textbf{Flaw:} The student fell into the \textbf{Incomplete Guess Trap}. When the forcing function is a polynomial of degree $n$, the guess must include all terms up to degree $n$, even if they are missing in $f(x)$. \\
\textbf{Correction:} The correct guess for $f(x) = x^2$ is a full quadratic: $y_p(x) = A x^2 + B x + C$.

\subsection*{Part 9: Proofs and Derivations}

\textbf{Solution 30:} \\
We want to show that $(D - \alpha)^2 [x e^{\alpha x}] = 0$. \\
First, apply $(D - \alpha)$ once. By the product rule, $D[x e^{\alpha x}] = e^{\alpha x} + \alpha x e^{\alpha x}$. \\
$(D - \alpha)[x e^{\alpha x}] = D[x e^{\alpha x}] - \alpha x e^{\alpha x} = (e^{\alpha x} + \alpha x e^{\alpha x}) - \alpha x e^{\alpha x} = e^{\alpha x}$. \\
Now apply $(D - \alpha)$ a second time to the result: \\
$(D - \alpha)[e^{\alpha x}] = D[e^{\alpha x}] - \alpha e^{\alpha x} = \alpha e^{\alpha x} - \alpha e^{\alpha x} = 0$. \\
Therefore, $(D - \alpha)^2 [x e^{\alpha x}] = 0$, proving that $(D - \alpha)^2$ is the annihilator.

\newpage

\section{Concept Checklist \& Question Mapping (Repeated)}
\label{sec:checklist_end}
\begin{itemize}
    \item \textbf{C1: Finding Differential Annihilators (Basic \& Composite)} - Questions: 1, 2, 3, 4, 5
    \item \textbf{C2: Determining the Form of $y_p$ (No Overlap)} - Questions: 6
    \item \textbf{C3: Determining the Form of $y_p$ (With Overlap)} - Questions: 7, 8, 9, 10
    \item \textbf{C4: Using the Annihilator Method to Determine $y_p$ Form} - Questions: 11, 12, 13, 14, 15
    \item \textbf{C5: Finding the General Solution (Full Solution)} - Questions: 16, 17, 18, 19, 20
    \item \textbf{C6: Initial Value Problems (IVPs)} - Questions: 21, 22, 23
    \item \textbf{C7: Systems of Differential Equations} - Questions: 24, 25
    \item \textbf{C8: Identifying Method Limitations} - Questions: 26, 27
    \item \textbf{C9: "Find the Flaw" / Conceptual Traps} - Questions: 28, 29
    \item \textbf{C10: Proofs and Derivations} - Questions: 30
\end{itemize}

\end{document}